\section{Held-Out Evaluation}
\label{sec:heldout}
\label{sec:expand}  % legacy label kept so appendix cross-references resolve

\emph{In one sentence:} on 119 fresh questions fully covered by four models, both six-vote ensembles separate from the single-sample expectation S1 with intervals excluding zero, the two sources remain statistically inseparable from each other, one fifth of the paraphrase headline is a single-call difficulty offset, and the residual margin is conditional --- it concentrates on quantitative tasks in a way that coincides with an instruction effect, and it reverses where paraphrases shift the task.

Every headline estimate in the development-suite analyses of \ref{app:dev} was computed on the same matrices that motivated the design choices, and per-model coverage there was truncated unevenly by free-tier quota events (33 to 90 questions per model). This section therefore carries the study's primary evaluation: a held-out suite that repeats the 16-call matrix on fresh questions, drawn by replaying the seeded shuffle before any of their outputs were inspected (Section~\ref{sec:method}). It takes the next unseen items from the same raw pools, with the same filters: 20 per benchmark, 19 for LSAT-LR where the filtered pool ran short, identifiers e000--e119 with zero overlap against the development suite. Paraphrases were generated by the same model under the same acceptance filters (embedding cosine at least 0.75, length ratio 0.4--2.5), so the pipeline differs from the main suite only in the identity of the drawn items, and item selection never consulted any model output. Coverage after the provider events of October 2026: Agnes3.0-Flash and Agnes2.5-Flash each completed all 119 questions; DeepSeekV4.1-Flash completed all 119 on a second free endpoint after its original provider exhausted its free tier at 9 questions, so its expansion row is a same-model provider switch rather than a mixed-provider composite (the 9-question original is retained only as a robustness record); DeepSeekV4-Pro-0813, a pro-tier sibling from the same family exposed by the second endpoint, completed all 119 questions there in direct-answer mode after the flash-tier sibling was moved behind a paid subscription mid-collection; Qwen3.7-Flash completed 60 of 119 before its provider revoked free-tier access, and because the runner drains multiple-choice items before GSM8K, the completed subset under-represents SAT-Math (1 of 20) and GSM8K (6 of 20); DeepSeekV4-Flash-0731 stopped at 13 questions and is excluded from tabulation. Table~\ref{tab:expand} reports the tabulated models, the pooled pair statistics over the four fully covered models, and the two rule replays. Because no greedy-paraphrase control exists on this suite (Section~\ref{sec:setup}), every paraphrase-ensemble margin below mixes prompt variation with per-call sampling noise, and we phrase the conclusions accordingly.

\begin{table}[tbp]
\centering
\footnotesize
\caption{Expansion set (119 fresh questions, 16-call matrix). Accuracy (\%) of the fixed conditions, the unanimity rate, and the two rule replays. $\S$: Qwen3.7-Flash completed 60 of 119 questions on a subset that under-represents SAT-Math and GSM8K (Section~\ref{sec:heldout}); its row is a partial replication, not a full one. Pooled aggregates the four fully covered models over 476 model--question pairs (questions weighted equally, each question contributing the mean over its four models).}
\label{tab:expand}
\scriptsize
\resizebox{\columnwidth}{!}{%
\begin{tabular}{l c c c c c c c c c c c c c}
\toprule
 & & \multicolumn{6}{c}{Accuracy (\%)} & & \multicolumn{2}{c}{Unanimity rule} & \multicolumn{3}{c}{Adaptive rule} \\
\cmidrule(lr){3-8}\cmidrule(lr){10-11}\cmidrule(lr){12-14}
Model & $n$ & Greedy & S1 & SC-6 & PE-6 & HB-6 & COMB & Unan.\% & Acc & $\bar{k}$ & Acc & $\bar{k}$ & Bud. \\
\midrule
Agnes3.0-Flash & 119 & 80.7 & 79.6 & 81.5 & 87.4 & 83.2 & 82.4 & 59 & 80.7 & 2.5 & 81.5 & 5.2 & 86\% \\
Agnes2.5-Flash & 119 & 89.9 & 88.1 & 90.8 & 91.6 & 89.1 & 89.9 & 76 & 90.8 & 2.3 & 89.9 & 4.7 & 79\% \\
DeepSeekV4.1-Flash & 119 & 73.1 & 77.6 & 79.8 & 81.5 & 83.2 & 83.2 & 58 & 80.7 & 2.8 & 84.0 & 5.2 & 87\% \\
DeepSeekV4-Pro-0813 & 119 & 79.0 & 78.0 & 80.7 & 82.4 & 82.4 & 84.0 & 61 & 80.7 & 2.7 & 83.2 & 5.1 & 85\% \\
Qwen3.7-Flash$^{\S}$ & 60 & 85.0 & 86.1 & 85.0 & 86.7 & 85.0 & 85.0 & 82 & 85.0 & 2.1 & 85.0 & 4.5 & 74\% \\
\midrule
Pooled & 476 & 80.7 & 80.8 & 83.2 & 85.7 & 84.5 & 84.9 & 63 & 83.2 & 2.6 & 84.7 & 5.1 & 84\% \\
\bottomrule
\end{tabular}}

\end{table}

\subsection{Primary contrasts: gains over S1}
\label{sec:heldoutprimary}

\emph{In one sentence:} every six- and twelve-vote condition separates from S1 with an interval excluding zero, and none of the pairwise orderings between conditions resolves.

On the fresh set, Agnes3.0-Flash shows the first ensemble gain in this study whose confidence interval excludes zero: PE-6 reaches 87.4\% against an S1 baseline of 79.6\%, a gain of $+7.8$ points (question-level bootstrap CI $+2.4$ to $+13.7$), while SC-6 gains $+2.0$ points (CI $-1.0$ to $+5.0$) and stays statistically indistinguishable from the single-call baseline --- exactly the main-suite pattern for sample-space ensembling. COMB-12 (82.4\%) lands between its two components, as before. Agnes2.5-Flash, now fully covered, replicates the direction on a higher and flatter baseline: PE-6 91.6\% against S1 88.1\% ($+3.5$, CI $+0.4$ to $+7.0$), SC-6 90.8\% ($+2.7$, CI $+0.7$ to $+4.9$), COMB-12 89.9\%. DeepSeekV4.1-Flash contributes the pattern in a different configuration: the pure strategies sit closer to the baseline (PE-6 81.5\% against S1 77.6\%, $+3.9$, CI $-0.8$ to $+8.7$; SC-6 79.8\%, $+2.2$, CI $-0.4$ to $+4.9$), and the two mixed allocations are the conditions whose intervals exclude zero for this model (HB-6 and COMB-12 both 83.2\%, $+5.6$ over S1, CIs $+1.5$ to $+9.8$ and $+1.7$ to $+9.7$). DeepSeekV4-Pro-0813 contributes a fourth configuration: its two individually significant conditions are the combined ensemble and, uniquely in the pool, pure self-consistency (COMB-12 84.0\%, $+6.0$ over S1, CI $+2.1$ to $+10.2$; SC-6 80.7\%, $+2.7$, CI $+0.1$ to $+5.5$), while the two paraphrase conditions fall between them with intervals spanning zero (PE-6 and HB-6 both 82.4\%, $+4.3$, CIs $-0.7$ to $+9.4$ and $-0.4$ to $+9.1$). Pooled over the four fully covered models (476 model--question pairs, cluster bootstrap on questions), every six- and twelve-vote condition separates from S1 with an interval excluding zero --- SC-6 $+2.4$ points (CI $+1.1$ to $+3.7$), PE-6 $+4.9$ points (CI $+2.3$ to $+7.6$), COMB-12 $+4.1$ points (CI $+2.0$ to $+6.3$) --- so on fresh questions both sources of diversity deliver measurable aggregation gains; designated as the study's primary contrasts (Section~\ref{sec:stats}), all three gains survive Holm adjustment within the four-contrast family (adjusted $p = 0.0008$, $<0.0001$, and $<0.0001$, respectively). The orderings between conditions, however, remain unresolvable at this sample size: PE$-$SC is $+2.5$ points (CI $-0.6$ to $+5.7$; Holm-adjusted $p = 0.14$), PE$-$HB $+1.3$ points (CI $-0.4$ to $+2.9$), and COMB$-$PE $-0.8$ points (CI $-2.9$ to $+1.3$), so no allocation separates from any other; in particular, the mixed allocations are among the best conditions for the two DeepSeek models, which limits how far a paraphrase-only recommendation can be pushed. The Qwen3.7-Flash partial replication behaves like the development-suite flagship models: 82\% of its questions are unanimous, every fixed condition lies within 1.1 points of S1, and both rule replays match SC-6 (85.0\%) at a fraction of the budget; its paraphrase-condition cells carry the self-preference caveat of Section~\ref{sec:setup}.

\subsection{The difficulty confound: what the $+4.9$ is made of}
\label{sec:heldoutds}

\emph{In one sentence:} paraphrases are slightly easier than original stems at a single call, so $+1.9$ of the $+4.9$ headline is a difficulty offset, and the like-for-like comparison is $+3.0$ within the paraphrase ensemble against $+2.4$ within resampling.

A difficulty confound must be separated from the ensembling effect before the margins can be read as aggregation: paraphrases could simply be easier than the original stems, and the comparison against S1 would then mix a difficulty margin with an aggregation margin. Table~\ref{tab:expand_ds} separates the two, and the pattern holds in all four models. On Agnes3.0 the paraphrases are only marginally easier at the single-call level (PE-1 81.2\% against S1 79.6\%, $+1.7$), so at most a fifth of its PE-6 margin is paraphrase difficulty, and the within-condition aggregation gains are $+6.2$ points inside the paraphrase ensemble (PE-6 87.4 against PE-1 81.2) against $+2.0$ inside resampling (SC-6 81.5 against S1 79.6). On Agnes2.5 the difficulty margin is similarly small (PE-1 89.2 against S1 88.1, $+1.1$), and the two aggregation gains converge ($\Delta$PE $+2.4$, $\Delta$SC $+2.7$). DeepSeekV4.1-Flash shows the same configuration at a smaller scale (PE-1 79.1 against S1 77.6, $+1.5$; $\Delta$PE $+2.4$, $\Delta$SC $+2.2$). The pro-tier DeepSeek shifts the balance: its paraphrases are the easiest in the pool at the single-call level (PE-1 81.1 against S1 78.0, $+3.1$), and its answer-space aggregation margin exceeds its paraphrase-side margin ($\Delta$SC $+2.7$ against $\Delta$PE $+1.3$), so its significant combined condition rests on the two sources jointly rather than on a dominant paraphrase component. Pooled over the four models, six-vote aggregation buys $+3.0$ points inside the paraphrase ensemble (PE-6 85.7 against PE-1 82.7) against $+2.4$ inside resampling (SC-6 83.2 against S1 80.8), with the single-call difficulty margin at $+1.9$ points contributing well under half of the pooled PE-6 advantage; the paraphrase-side aggregation margin is roughly 1.3 times the answer-space margin, and because $+1.9$ of the $+4.9$ headline gap is a difficulty offset, the aggregation-against-aggregation comparison is $+3.0$ versus $+2.4$ rather than the raw four-versus-two spread. All of these margins mix prompt variation with per-call sampling noise (Section~\ref{sec:conditions}); no noise-free paraphrase control exists on this suite.

\begin{table}[tbp]
\centering
\footnotesize
\caption{Expansion-set decomposition for the fully covered models: single-call accuracy of original-stem samples (S1) and paraphrase first-samples (PE-1) against the six-vote ensembles (SC-6, PE-6), with the within-condition ensembling gains ($\Delta$SC $=$ SC-6 $-$ S1, $\Delta$PE $=$ PE-6 $-$ PE-1), by dataset. Pooled rows weight questions equally.}
\label{tab:expand_ds}
\scriptsize
\resizebox{\columnwidth}{!}{%
\begin{tabular}{l l c c c c c c c c}
\toprule
 & & & \multicolumn{2}{c}{Single call (\%)} & & \multicolumn{2}{c}{6-vote (\%)} & \multicolumn{2}{c}{Ensembling gain (pp)} \\
\cmidrule(lr){4-5}\cmidrule(lr){6-7}\cmidrule(lr){8-9}
Model & Dataset & $n$ & S1 & PE-1 & SC-6 & PE-6 & $\Delta$SC & $\Delta$PE \\
\midrule
Agnes3.0-Flash & gsm8k & 20 & 95.0 & 97.5 & 95.0 & 100.0 & +0.0 & +2.5 \\
Agnes3.0-Flash & logiqa & 20 & 75.0 & 63.3 & 90.0 & 65.0 & +15.0 & +1.7 \\
Agnes3.0-Flash & lsat\_lr & 19 & 86.8 & 88.6 & 89.5 & 94.7 & +2.6 & +6.1 \\
Agnes3.0-Flash & piqa & 20 & 88.3 & 86.7 & 90.0 & 95.0 & +1.7 & +8.3 \\
Agnes3.0-Flash & sat\_math & 20 & 61.7 & 82.5 & 60.0 & 100.0 & -1.7 & +17.5 \\
Agnes3.0-Flash & truthfulqa & 20 & 70.8 & 69.2 & 65.0 & 70.0 & -5.8 & +0.8 \\
 & pooled & 119 & 79.6 & 81.2 & 81.5 & 87.4 & +2.0 & +6.2 \\
\midrule
Agnes2.5-Flash & gsm8k & 20 & 94.2 & 95.8 & 95.0 & 100.0 & +0.8 & +4.2 \\
Agnes2.5-Flash & logiqa & 20 & 73.3 & 77.5 & 70.0 & 80.0 & -3.3 & +2.5 \\
Agnes2.5-Flash & lsat\_lr & 19 & 90.4 & 94.7 & 100.0 & 94.7 & +9.6 & +0.0 \\
Agnes2.5-Flash & piqa & 20 & 94.2 & 93.3 & 100.0 & 95.0 & +5.8 & +1.7 \\
Agnes2.5-Flash & sat\_math & 20 & 98.3 & 99.2 & 100.0 & 100.0 & +1.7 & +0.8 \\
Agnes2.5-Flash & truthfulqa & 20 & 78.3 & 75.0 & 80.0 & 80.0 & +1.7 & +5.0 \\
 & pooled & 119 & 88.1 & 89.2 & 90.8 & 91.6 & +2.7 & +2.4 \\
\midrule
DeepSeekV4.1-Flash & gsm8k & 20 & 98.3 & 98.3 & 100.0 & 100.0 & +1.7 & +1.7 \\
DeepSeekV4.1-Flash & logiqa & 20 & 52.5 & 65.0 & 60.0 & 65.0 & +7.5 & +0.0 \\
DeepSeekV4.1-Flash & lsat\_lr & 19 & 79.8 & 84.2 & 78.9 & 84.2 & -0.9 & +0.0 \\
DeepSeekV4.1-Flash & piqa & 20 & 79.2 & 72.5 & 80.0 & 75.0 & +0.8 & +2.5 \\
DeepSeekV4.1-Flash & sat\_math & 20 & 75.8 & 74.2 & 75.0 & 80.0 & -0.8 & +5.8 \\
DeepSeekV4.1-Flash & truthfulqa & 20 & 80.0 & 80.8 & 85.0 & 85.0 & +5.0 & +4.2 \\
 & pooled & 119 & 77.6 & 79.1 & 79.8 & 81.5 & +2.2 & +2.4 \\
\midrule
DeepSeekV4-Pro-0813 & gsm8k & 20 & 94.2 & 95.0 & 95.0 & 95.0 & +0.8 & +0.0 \\
DeepSeekV4-Pro-0813 & logiqa & 20 & 54.2 & 68.3 & 55.0 & 70.0 & +0.8 & +1.7 \\
DeepSeekV4-Pro-0813 & lsat\_lr & 19 & 78.1 & 83.3 & 84.2 & 84.2 & +6.1 & +0.9 \\
DeepSeekV4-Pro-0813 & piqa & 20 & 90.0 & 83.3 & 95.0 & 80.0 & +5.0 & -3.3 \\
DeepSeekV4-Pro-0813 & sat\_math & 20 & 69.2 & 79.2 & 70.0 & 85.0 & +0.8 & +5.8 \\
DeepSeekV4-Pro-0813 & truthfulqa & 20 & 82.5 & 77.5 & 85.0 & 80.0 & +2.5 & +2.5 \\
 & pooled & 119 & 78.0 & 81.1 & 80.7 & 82.4 & +2.7 & +1.3 \\
\midrule
Pooled & gsm8k & 80 & 95.4 & 96.7 & 96.2 & 98.8 & +0.8 & +2.1 \\
Pooled & logiqa & 80 & 63.8 & 68.5 & 68.8 & 70.0 & +5.0 & +1.5 \\
Pooled & lsat\_lr & 76 & 83.8 & 87.7 & 88.2 & 89.5 & +4.4 & +1.8 \\
Pooled & piqa & 80 & 87.9 & 84.0 & 91.2 & 86.2 & +3.3 & +2.3 \\
Pooled & sat\_math & 80 & 76.2 & 83.8 & 76.2 & 91.2 & +0.0 & +7.5 \\
Pooled & truthfulqa & 80 & 77.9 & 75.6 & 78.8 & 78.8 & +0.8 & +3.1 \\
Pooled & pooled & 476 & 80.8 & 82.7 & 83.2 & 85.7 & +2.4 & +3.0 \\
\bottomrule
\end{tabular}}

\end{table}

A cross-suite consistency check pools every question completed by the two models covered in both suites (development plus held-out: 152 questions for Agnes3.0-Flash, 159 for Agnes2.5-Flash). The held-out pattern survives the union: on Agnes3.0-Flash, PE-6 sits 6.5 points above S1 (CI $+2.2$ to $+11.1$) while SC-6 gains $+1.9$ (CI $-0.7$ to $+4.4$); on Agnes2.5-Flash the two sources converge ($+2.2$ each, with the SC-6 interval at $+0.6$ to $+3.9$ and the PE-6 interval at $-0.4$ to $+4.9$). The per-model point estimates nonetheless move several points between suites (Agnes2.5-Flash: PE-6 minus S1 is $-1.7$ on its 40-question development subset, $+3.5$ on the 119 held-out questions, and $+2.2$ on the 159-question union; the within-condition $\Delta$PE uses PE-1 and moves in the same direction), a spread that is what unevenly truncated small samples deliver and that cautions against reading any single per-model cell as an effect size; the pooled four-model contrasts, not per-model cells, are where the study's power lies.

\subsection{Dataset-level mechanism, and the instruction effect}
\label{sec:heldoutmech}

\emph{In one sentence:} the paraphrase-side margin concentrates on quantitative tasks, and an instruction-variant analysis on the same logs shows this concentration coincides with an instruction effect --- the canonical directive under-elicits, the best single variant (selected post hoc) already exceeds the six-variant ensemble, and the ensemble captures most but not all of that effect.

The dataset decomposition carries the mechanism, with two cautions stated up front: each cell rests on twenty questions per model (nineteen for LSAT-LR), and all gains below are the within-condition margins of Table~\ref{tab:expand_ds} ($\Delta$SC $=$ SC-6 $-$ S1, $\Delta$PE $=$ PE-6 $-$ PE-1), with single-call difficulty margins reported alongside wherever they are large. On Agnes3.0-Flash's SAT-Math subset, six samples of the original stem move accuracy from 61.7\% to 60.0\% ($\Delta$SC $-1.7$) --- resampling re-derives the same computations --- while the paraphrase ensemble recovers the subset almost entirely (PE-6 100.0\%, $\Delta$PE $+17.5$); all eight SAT-Math questions on which SC-6 and PE-6 disagree are won by the paraphrase vote. Because SAT-Math paraphrases vary the instruction wording rather than the stem, part of this cell is an instruction effect rather than an aggregation effect: the paraphrase single call already sits 20.8 points above S1 there (PE-1 82.5 against 61.7), which shows that the canonical directive to answer without explanation is itself suboptimal for quantitative items and that some instruction variants elicit more careful behaviour. The within-condition margin is net of this base-rate shift, but the mechanism reading must be stated carefully, and the instruction-variant analysis makes it precise: scoring each of the six handcrafted instruction variants separately, the best variant --- selected post hoc across variants on the same data, hence an optimistic ceiling --- lifts single-call accuracy on the four models' pooled SAT-Math items from 76.2\% under the canonical directive to 95.0\% (Agnes3.0-Flash 61.7 to 100.0, DeepSeekV4.1-Flash 75.8 to 90.0, the pro tier 69.2 to 90.0, with Agnes2.5-Flash already at ceiling), whereas on PIQA, whose instruction variants reword the directive without changing elicitation quality, the best variant sits at the canonical level (86.2\% against 87.9\% pooled). On this data the instruction effect is present, specific to quantitative items, and larger pooled than the aggregation margins themselves --- a post-hoc reading, since the variant was selected on the same items; the six-variant paraphrase ensemble, which requires no variant selection, captures about eighty percent of it (pooled PE-6 91.2\% against a 76.2\% canonical single call and the 95.0\% post-hoc ceiling, and it reaches the ceiling exactly on Agnes3.0-Flash). What these logs cannot decide is whether six resamples under the single best instruction would aggregate a further margin on top of the best variant's single call; that condition is absent from the matrix by design, and we flag it explicitly as the one unrun control of this cell\TODO{pre-registered follow-up: SC-6 under the pre-selected best instruction variant on a fresh SAT-Math draw --- script provided with the submission, result pending}. A pre-registration on the development logs has already been fixed: scoring the six variants on the development-suite SAT-Math items selects variant P1 at 89.9\% pooled single-call accuracy, slightly \emph{below} the canonical directive there (90.6\%), so on the development draw the canonical instruction was not suboptimal --- the fresh-draw test, not the post-hoc ceiling, is what can decide the question. GSM8K shows the same direction in weaker form (PE-6 100.0\%, $\Delta$PE $+2.5$ on a floor of 95.0\%). LogiQA is the mirror image for this model: paraphrases are 11.7 points harder at the single-call level (PE-1 63.3 against S1 75.0), indicating that restating a verbal-reasoning stem can shift the reasoning problem itself, and there the sample vote is the only condition that aggregates a gain (SC-6 90.0\%, $\Delta$SC $+15.0$ against $\Delta$PE $+1.7$; five of the seven flips favour SC). TruthfulQA is the one dataset where its six-sample vote actively hurts (SC-6 65.0 against S1 70.8), consistent with the development-suite reading that on knowledge questions the modal answer can be a confident misconception; on this dataset the paraphrase side retains a small within-condition margin ($\Delta$PE $+3.1$ pooled over the four models; on Agnes3.0 alone it is $+0.8$), which we report as a residual prompt-side effect on a mixed-locus dataset rather than as a neutral reading.

Agnes2.5-Flash contributes the counter-case that bounds the mechanism: its SAT-Math and GSM8K floors sit near ceiling (S1 98.3 and 94.2), its pooled paraphrase-side aggregation margin is correspondingly small ($\Delta$PE $+2.4$), and on LogiQA the Agnes3.0 ordering reverses (PE-6 80.0 against SC-6 70.0, $\Delta$PE $+2.5$ against $\Delta$SC $-3.3$), so on stem-locus datasets the PE--SC ordering is a property of the model--paraphrase interaction rather than of the dataset. The DeepSeek models sharpen the same caution with opposite difficulty signs: DeepSeekV4.1's LogiQA paraphrases are 12.5 points easier at the single-call level (PE-1 65.0 against S1 52.5) --- where Agnes3.0's were 11.7 points harder --- and there the sample vote is again the only condition that aggregates a gain ($\Delta$SC $+7.5$ against $\Delta$PE $+0.0$), while its SAT-Math repeats the quantitative direction ($\Delta$SC $-0.8$, $\Delta$PE $+5.8$) and its TruthfulQA supplies the counter-case to Agnes3.0's, with the sample vote helping ($\Delta$SC $+5.0$). The pro-tier DeepSeek shows the pool's largest single-call difficulty margins in both directions. Its LogiQA paraphrases are 14.1 points easier (PE-1 68.3 against S1 54.2); the paraphrase ensemble holds that base-rate advantage through voting (PE-6 70.0 against SC-6 55.0), but the within-condition gains are small ($\Delta$PE $+1.7$ against $\Delta$SC $+0.8$), so what moves this cell is the difficulty margin rather than aggregation --- the mirror image of Agnes3.0's LogiQA. On PIQA its paraphrases are 6.7 points harder and the paraphrase ensemble loses to the sample ensemble by fifteen points (PE-6 80.0 against SC-6 95.0; measured against S1 the paraphrase side loses 10.0 points, of which 6.7 is the harder single-call base and $-3.3$ the within-condition change, while the sample side gains $\Delta$SC $+5.0$). Pooled over the four models, the quantitative datasets carry the paraphrase-side aggregation advantage (SAT-Math $\Delta$PE $+7.5$ against $\Delta$SC $+0.0$; GSM8K $+2.1$ against $+0.8$), while on LogiQA the sample side aggregates the larger gain ($\Delta$SC $+5.0$ against $\Delta$PE $+1.5$), the paraphrases' pooled single-call advantage there ($+4.8$ over S1) doing more of the work than aggregation; TruthfulQA is near-neutral for resampling ($\Delta$SC $+0.8$) while retaining a small paraphrase-side margin ($\Delta$PE $+3.1$), so the sample vote neither gains nor loses much there while the paraphrase vote keeps a modest advantage on this mixed-locus dataset. Two readings of these cells deserve explicit caution. First, the single-call difficulty margins of the same dataset's paraphrases swing by more than twelve points across models (LogiQA: $-11.7$ on Agnes3.0, $+12.5$ on DeepSeekV4.1, $+14.1$ on the pro tier) --- too large to dismiss as sampling noise at these denominators, and consistent either with specific rewrites shifting the task for specific models or with equivalence failures that the aggregate validation filters cannot see; the model--paraphrase interaction is real but not currently predictable, and we treat it as descriptive. Second, because each cell is a twenty-question stratum, individual margins carry wide intervals; the transferable content is the pair of directions --- paraphrase-side aggregation pays on quantitative tasks, in a way that coincides with the instruction effect, and reverses where paraphrases shift the task --- rather than any single cell's magnitude. Both the instruction-variant selection (post hoc, on the same data it scores) and the absence of the best-instruction SC-6 control are stated as limitations in Section~\ref{sec:discussion}.

\subsection{Strata decomposition}
\label{sec:heldoutstrata}

\emph{In one sentence:} the (statistically unresolved) PE$-$SC margin concentrates in the weak-majority stratum, exactly where answer-space votes have the least to aggregate, and part of that concentration is mechanical.

The strata defined on the development suite also organize the held-out votes, and the decomposition --- computed on the same twelve votes that define the pooled contrasts --- locates where the (statistically unresolved) PE$-$SC margin concentrates. The weak-majority stratum (mode share below two-thirds) holds 24, 8, 14, and 15 questions for the four models: PE-6 scores 58.3\% against 33.3\% for SC-6 and 31.9\% for S1 on Agnes3.0, 50.0\% against 25.0\% and 27.1\% on Agnes2.5, 50.0\% against 35.7\% and 33.3\% on DeepSeekV4.1, and 53.3\% against 33.3\% and 31.1\% on the pro tier; averaging the four models' strata (61 pairs in total) gives PE-6 52.9\% against 31.8\% for SC-6 and 30.9\% for S1, a $+22$-point margin exactly where answer-space votes have the least to aggregate, and for the two DeepSeek models the combined twelve-vote ensemble adds a further margin on top (COMB-12 57.1\% against PE-6 50.0\% on V4.1; 60.0\% against 53.3\% on the pro tier). We emphasize what this decomposition is and is not: it reuses the same votes that produce the pooled PE$-$SC gap of $+2.5$ points (interval touching zero), so it is a localization of that gap, not an independent replication of it, and part of the concentration is mechanical --- a question on which the sample side is degenerate but wrong is exactly a question where only the paraphrase side can move the answer, so the stratum where PE can win is selected, in part, by the failure of SC. The unanimous stratum ($n=70$, $n=90$, $n=69$, and $n=73$) is saturated at 96.3\% pooled under every condition, and the strong-majority stratum ($n=25$, $n=21$, $n=36$, and $n=31$) sits between, with the two six-vote ensembles at 76--77\% against an S1 of 67.0\%.

\subsection{Greedy versus S1 on the held-out suite}
\label{sec:heldoutgreedy}

\emph{In one sentence:} the greedy--S1 divergence reappears on exactly one of the four held-out endpoints, so it is a property of specific model--endpoint configurations, not a general property of greedy decoding.

The $\tau=0$ anomaly appears on exactly one of the four models, and the split is informative. On the two Agnes endpoints S1 minus greedy is $-1.1$ and $-1.8$ points, and on the pro-tier DeepSeek $-1.0$. On the direct-answer DeepSeekV4.1 endpoint the development-suite signature reappears: greedy sits 4.5 points below S1 (73.1\% against 77.6\%), and auditing its 16 off-mode greedy answers shows 13 of them wrong, with completion lengths of one to eleven characters --- genuine divergence between the temperature-zero answer and the sampled mode, not the truncation artifact of the development suite; the pro-tier sibling's 14 off-mode answers split 10 wrong against 4 right, a third clean case. Pooled, 50 of 476 greedy answers (10.5\%) fall outside the six-sample mode. This relocates the greedy audit's lesson: the temperature-zero failure is a property of the model--endpoint configuration rather than of the question draw or of a universal decoding pathology --- absent on three of four held-out endpoints, present without truncation on the fourth --- so practitioners should verify greedy against the sampled expectation per deployment rather than assume either baseline.

\subsection{Fragility on the held-out suite}
\label{sec:heldoutfragility}

\emph{In one sentence:} fragility persists at a higher rate than the truncated development pool suggested and loses its clean locus split --- instruction-locus flips do occur on fresh questions, almost entirely on SAT-Math.

Because the fragility analysis was established on the development suite, the held-out matrices recompute it, and the reading changes in ways worth stating. Fragility persists at a higher rate than the truncated development pool suggested, and it loses its clean locus split. Pooled over the four fully covered models, 17.4\% of stem-locus questions flip their majority answer under paraphrasing (41 of 236 model--question pairs), with LogiQA again carrying the concentration (40, 15, 35, and 35\% per model), while instruction-locus questions now flip as well, at 9.6\% pooled (23 of 240 pairs), where the development suite saw zero of 52. The instruction-locus flips sit almost entirely on SAT-Math (40, 0, 15, and 20\% per model, against at most 15\% on PIQA and 5\% on GSM8K), which is the dataset whose instruction variants differ in how well they elicit careful quantitative answering; the development suite's zero was a property of its question draw and of its near-ceiling instruction-locus items, not a structural law. What survives on the fresh questions is the direction of the asymmetry --- stem rewordings flip roughly twice as often pooled --- and the mechanism: flips concentrate in the model--task cells that carry answer entropy and ensembling gains.

\subsection{Error signals at adequate error counts}
\label{sec:heldouterror}

\emph{In one sentence:} at twelve to twenty-one errors per model both agreement scores carry signal, their coverages overlap only partially, and the summed disagreement score is numerically the strongest pooled detector --- a descriptive ranking whose interval overlaps the singles'.

\begin{table}[tbp]
\centering
\footnotesize
\caption{Held-out error-signal analysis. Reference errors are those of the COMB-12 ensemble; detectors are reversed agreement scores ($c_S$ over the six canonical samples, $c_P$ over the six paraphrase first decodes, Both their summed disagreement), reported as AUROC with question-level bootstrap 95\% intervals (cluster bootstrap over questions for the pooled row); every fully covered model carries twelve to twenty-one reference errors, so no row is descriptive-only. Flagged counts the reference errors on which the given source shows any dissent; by construction the union of the two flagged sets contains each single set, so Either-coverage is an upper bound by definition. $^{\S}$Qwen3.7-Flash partial replication (60 of 119 questions) with the self-preference caveat of Section~\ref{sec:setup}; excluded from the pooled row.}
\label{tab:uncertexp}
\footnotesize
\begin{tabular*}{\columnwidth}{@{\extracolsep{\fill}} l c c c c c c c}
\toprule
 & & \multicolumn{3}{c}{AUROC [95\% CI]} & \multicolumn{3}{c}{Flagged errors} \\
\cmidrule(lr){3-5}\cmidrule(lr){6-8}
Model & Errors & $c_S$ & $c_P$ & Both & $c_S$ & $c_P$ & Either \\
\midrule
Agnes3.0-Flash & 21 & 0.80 [0.69, 0.90] & 0.81 [0.71, 0.90] & 0.86 [0.78, 0.93] & 16/21 & 18/21 & 20/21 \\
Agnes2.5-Flash & 12 & 0.74 [0.58, 0.89] & 0.84 [0.68, 0.98] & 0.83 [0.67, 0.97] & 7/12 & 9/12 & 9/12 \\
DeepSeekV4.1-Flash & 20 & 0.70 [0.58, 0.82] & 0.62 [0.50, 0.75] & 0.73 [0.61, 0.84] & 13/20 & 9/20 & 16/20 \\
DeepSeekV4-Pro-0813 & 19 & 0.74 [0.61, 0.86] & 0.69 [0.56, 0.82] & 0.80 [0.69, 0.90] & 13/19 & 10/19 & 16/19 \\
Qwen3.7-Flash$^{\S}$ & 9 & 0.72 [0.56, 0.90] & 0.78 [0.60, 0.94] & 0.80 [0.61, 0.97] & 4/9 & 6/9 & 6/9 \\
\midrule
Pooled (476 pairs) & 72 & 0.75 [0.67, 0.84] & 0.74 [0.65, 0.83] & 0.81 [0.73, 0.89] & 49/72 & 46/72 & 61/72 \\
\bottomrule
\end{tabular*}

\end{table}

Table~\ref{tab:uncertexp} reports the error-signal analysis at error counts the development models could not supply, so every point estimate now carries an interval instead of the descriptive-only dagger of \ref{app:dev}. Four readings emerge, of which two are corrections of development-suite readings. First, both signals are validated as ensemble-error detectors: pooled, sample agreement reaches an AUROC of 0.75 (95\% CI 0.67 to 0.84) and paraphrase agreement 0.74 (0.65 to 0.83). Second, the development-suite claim of a per-model dominant signal does not survive: on the development suite, sample agreement had looked decisively stronger than paraphrase agreement for Agnes3.0-Flash (0.96 against 0.60); that reading rested on three errors --- below the five-error rule of Section~\ref{sec:method} --- and the held-out recomputation at eighteen errors overturns it: the two signals are comparable there (paraphrase 0.81 against sample 0.80), and only Agnes2.5-Flash shows a clear paraphrase advantage (0.84 against 0.74), while the two DeepSeek tiers lean the other way (0.70 and 0.74 for sample against 0.62 and 0.69 for paraphrase). The stronger single signal is therefore model-dependent, and the earlier apparent regularity --- ``the stronger signal tracks the stronger diversity source'' --- is not supported at adequate error counts. Third, the coverage complementarity the development suite suggested holds at scale: of the 72 pooled errors, sample dissent flags 49 and paraphrase dissent 46, fifteen are flagged only by the sample side and twelve only by the paraphrase side, so either probe alone misses 23 to 26 errors where the pair misses eleven. The union's flagged set contains each single set by construction, so the pair's coverage is an upper bound by definition; the empirical content is the partial disjointness --- fifteen and twelve errors reachable only through one source --- not the union size. Fourth, the sum of the two disagreement scores --- which the development suite could not endorse --- is numerically the strongest pooled detector on the fresh questions (0.81, CI 0.73 to 0.89); we label this ranking descriptive, because its interval overlaps those of the single signals (0.75 and 0.74), and on the per-model rows it is the best or within noise of the best single signal everywhere rather than uniformly ahead. The structural caveat of Section~\ref{sec:method} still bounds how far these numbers should be read: on the full question set the AUROC partly reflects the definitional degenerate stratum, where a unanimous but wrong reference is unflaggable by construction; questions with any dissent among the twelve votes carry a 35.1\% error rate pooled against 3.6\% on unanimous questions, and within the non-degenerate stratum alone the per-model AUROCs fall to between 0.49 and 0.86. Agreement scores, in other words, identify which questions sit in the risky regime far better than they rank risk among those questions --- the same regime-first reading that the accuracy analysis delivers.

\subsection{Transfer of the adaptive allocation rule}
\label{sec:heldoutrule}

\emph{In one sentence:} the rule saves 13 to 21 percent of the calls on fresh questions, exceeds every fixed allocation on one model by a single question, and loses one to seven questions to the best fixed conditions elsewhere.

The adaptive allocation rule was designed on the development-suite matrices (its development-suite simulation is reported in \ref{app:dev}), so the held-out set supplies the fresh-question transfer test. On three of the four fully covered models the rule transfers in budget but not in allocation. It reproduces SC-6 exactly (no question gap) at 86\% of the six-call budget on Agnes3.0 (81.5\% at a mean of 5.2 calls), and it matches SC-6 exactly at 74\% of budget on the Qwen partial (85.0\%). But it does not capture the paraphrase-side margin there: it sits seven questions below PE-6 on Agnes3.0, because questions whose two-call probe shows diversity on both sides are routed to the hybrid vote, which underperforms a pure paraphrase vote on this distribution. On Agnes2.5 the rule lands one question below SC-6 (89.9\% against 90.8\%) and two below PE-6, at 79\% of budget (a mean of 4.7 calls); its probe routes 97 of 119 questions to the four-call unanimous exit, 10 to the sample vote, 9 to the paraphrase vote, and 3 to the hybrid. On DeepSeekV4.1 the rule does what no fixed condition does there: 84.0\% at a mean of 5.2 calls (87\% of budget), five questions above SC-6, three above PE-6, and one above HB-6 and COMB-12, the only model in either suite on which the rule exceeds every fixed allocation --- though the margin over the best of them is a single question, so we read it as reaching the fixed frontier rather than dominating it; its probe routes 83 of 119 questions to the unanimous exit, 18 to the paraphrase vote, 12 to the sample vote, and 6 to the hybrid. On the pro-tier DeepSeek the rule lands at 83.2\% at a mean of 5.1 calls (85\% of budget), three questions above SC-6 and one above PE-6 but one question below its best fixed condition, the combined ensemble (84.0\%); its probe routes 86 of 119 questions to the unanimous exit, 20 to the paraphrase vote, 9 to the sample vote, and 4 to the hybrid. Pooled over the four models the rule returns 84.7\% at a mean of 5.1 calls (84\% of budget). The development-suite conclusion therefore needs its held-out amendment: the rule saves 13 to 21 percent of the calls, exceeds every fixed allocation on one model (by a single question), loses one question to its best fixed condition on the pro tier and two to PE-6 on Agnes2.5, and sits seven questions below the paraphrase vote on Agnes3.0 --- routing toward paraphrases when answer-space diversity is low remains the regime where the most is gained, and it is exactly the regime the probe's two-call answer-space check undersamples.

\subsection{What the two suites jointly establish}
\label{sec:heldoutjoint}

\emph{In one sentence:} six-call resampling moves accuracy little for any model in either suite, the paraphrase ensemble produces the largest well-powered gains, and both statements hold only under the two qualifying conditions --- quantitative tasks whose canonical instruction under-elicits, and paraphrases that preserve task semantics.

These results qualify, without overturning, the development-suite null. Read together, the two suites say that six-call sample-space ensembling moves accuracy little for any model in either suite (per-model $\Delta$SC between $-3.3$ and $+2.7$ points), while the paraphrase ensemble --- each of whose votes mixes prompt variation with decoding noise (Section~\ref{sec:conditions}) --- produces the largest well-powered gains: pooled $+4.9$ points over S1 on the held-out set (CI $+2.3$ to $+7.6$), of which $+1.9$ is a single-call difficulty offset and $+3.0$ the within-condition aggregation gain, against $+2.4$ for self-consistency (CI $+1.1$ to $+3.7$). The margin is conditional in a way that the title question's answer must reflect: the paraphrase-ensemble margin pays where the canonical instruction under-elicits --- on quantitative tasks, where the best single instruction variant (selected post hoc) already exceeds the six-variant ensemble and the ensemble captures most, not all, of that instruction effect --- and it reverses where paraphrases shift the task's verbal semantics. The allocation that realizes the gain is model-dependent --- the pure paraphrase vote on the two Agnes models, the hybrid and combined allocations on DeepSeekV4.1, the combined ensemble and self-consistency on its pro-tier sibling --- so the transferable content is the two qualifying conditions rather than a recipe. The gain is established here for four fully covered models from two providers and two families (two Agnes generations, two DeepSeek tiers) on one fresh draw, so we read it as a demonstrated possibility with an identified mechanism rather than as a universal effect.
